\begin{abstract}
Este trabajo presenta una implementación computacional del modelo teórico de David Anisi sobre el mercado laboral, utilizando técnicas de Modelización Basada en Agentes (ABM). El modelo reproduce la dinámica del empleo, el trabajo doméstico y los índices de bienestar laboral a través de la interacción heterogénea entre tres tipos de agentes: Administración Pública, Empresas y Trabajadores. 

Nuestra contribución principal es triple: (1) la traducción del modelo matemático de Anisi a un sistema de agentes interactivos implementado en Python con la librería Mesa, (2) la calibración del modelo para el caso específico de España utilizando datos macroeconómicos de 2024, y (3) el desarrollo de una interfaz web interactiva que permite la simulación en tiempo real y la visualización dinámica de los indicadores clave del mercado laboral.

Los resultados demuestran que el modelo captura adecuadamente las dinámicas fundamentales descritas por Anisi, incluyendo la tasa de paro, el índice de colocación (I1), el índice de dualismo (I2) y el índice de ocupación/frustración (I3). La herramienta de visualización desarrollada permite explorar el impacto de políticas públicas como cambios en el gasto público, la presión fiscal y la regulación de la jornada laboral.

\textbf{Palabras clave:} Modelización Basada en Agentes, Mercado Laboral, Modelo de Anisi, Simulación Económica, España
\end{abstract}
